%% kapa: exam=1 round=27 year=2016 period=2 form=A subject=appraisal_law subject_ko=감정평가관계법규 seq=40 number=40
%% source: Q-Net 감정평가사 기출문제 (cst00309) artlSeq=5169402 file=extracted/1008/5169402/제27회 감정평가사 1차 문제/제27회 감정평가사 2교시문제(A형).pdf
%% conversion: kordoc 4.24.51 -> tools/gen1.py (auto); review=proofread-sonnet; flags=none
\begin{question}{40}
도시 및 주거환경정비법령상 비용부담 등에 관한 설명으로 옳지 않은 것은?
\begin{choices}
  \item 시장ㆍ군수는 시장ㆍ군수가 아닌 사업시행자가 시행하는 정비사업의 정비계획에 따라 설치되는 도시ㆍ군계획시설 중 녹지에 대하여는 그 건설에 소요되는 비용의 전부 또는 일부를 부담할 수 있다.
  \item 시장ㆍ군수는 그가 시행하는 정비사업으로 인하여 현저한 이익을 받는 정비기반시설의 관리자가 있는 경우에는 그 정비기반시설의 관리자와 협의하여 당해 정비사업비의 3분의 2까지를 그 관리자에게 부담시킬 수 있다.
  \item 시장ㆍ군수가 아닌 사업시행자는 부과금 또는 연체료를 체납하는 자가 있는 때에는 시장ㆍ군수에게 그 부과ㆍ징수를 위탁할 수 있다.
  \item 공동구에 수용될 전기ㆍ가스ㆍ수도의 공급시설과 전기통신시설 등의 관리자가 부담할 공동구의 설치에 소요되는 비용의 부담비율은 공동구의 점용예정면적비율에 의한다.
  \item 사업시행자가 정비사업비의 일부를 정비기반시설의 관리자에게 부담시키고자 하는 때에는 정비사업에 소요된 비용의 명세와 부담 금액을 명시하여 그 비용을 부담시키고자 하는 자에게 통지하여야 한다.
\end{choices}
\end{question}
\answer{2}
